\section{Results II: negative results as a taxonomy}
\label{sec:negatives}

Fourteen of the original build's 25 registered predictions failed as frozen. % F-147, F-148 (C-22)
Appendix~\ref{app:predictions} lists every one with its verdict and artifact; the frozen bars are in the public design documents and the measured values in each artifact's results file; none required changing any frozen parameter to diagnose. Here we do two things: classify the failures, and work one example of each class in enough detail that the reader can judge whether the diagnostic readings are honest. The classification (\gInterp): % C-23

\begin{enumerate}
\item \textbf{Bounds written at the wrong scale.} The prediction constrained the right quantity, but its bar was set at the scale of a different one.
\item \textbf{Thresholds frozen too strictly.} The predicted effect exists and points the right way, but the frozen numeric bar was set beyond what the effect delivers.
\item \textbf{Structure the prediction mischaracterized.} The measurement is informative and the machinery correct, but the world disagreed with the registered hypothesis.
\end{enumerate}

\paragraph{Worked example, class 1 (P1.1).} We registered that the full-schedule residual trace would fall below $10^{-2}$ at $p = 1/20$ and shrink $10\times$ from distance 3 to 5. Measured: $0.0837$ and $0.0740$, a ratio of $1.13$. % F-061, F-106
The bound was written at logical-error-probability scale, but $\operatorname{tr}\Xi$ is the linear-estimator MMSE of a $\pm1$ observable---a different quantity with a different scale and different distance behavior. The failure is real and taught us something concrete: the residual's units are variance of an observable, not failure probability, and predictions must be written in the object's own units. The linear-class limitation the number reflects is independently visible error-by-error: an exact bounded-weight cross-check shows that the affine estimator $A^{*}$ built from REP(5)'s four syndrome bits reproduces majority-logic minimum-weight decoding on 14 of 16 low-weight errors and fails on two (\gExact). That the failure is a limitation of the whole linear class, not of this one estimator, is our reading (\gInterp). % F-043

\paragraph{Worked example, class 2 (P6.2).} We registered that a slack point exists strictly before budget exhaustion, at tolerance $\lambda_{\mathrm{tol}} = 10^{-9}$. Measured: $b^{*} = b_{\max}$ for both REP(5) and SURF(3), because even the smallest genuine marginal value in these families exceeds the frozen tolerance. % F-072, F-058
The value curve itself is exact and its tail marginals are tiny (of order $10^{-5}$ to $10^{-4}$), % F-057
so the intended phenomenon---most of the budget buys almost nothing---is visibly present in the curve; the certificate failed because its threshold was frozen four to five orders of magnitude below the phenomenon. % F-057, F-058
A defensible re-thresholding would pass, and we did not do it, because a moved threshold after seeing data is exactly what the discipline forbids.

\paragraph{Worked example, class 3 (P5.3).} The protocol-trap prediction of Section~\ref{sec:results-memory}: we registered that our alternating schedule was an artifact that internalization would dissolve. Internalization preserved all 4 witnesses; the schedule carries genuine predictive content. % F-070
The machinery did precisely its job---it is the registered hypothesis that was wrong, and the failure sharpened the concept: ``externally scheduled'' does not imply ``predictively empty.'' For this package and its declared observation catalog, internalizing the schedule left every witness intact: the phase carried information about future observables, so the trap did not classify it as an artifact.

One more failure deserves mention here because the entire second half of the paper grows out of it: P2.1, the deployable drift witness that never detected within the tested grid (class 3, with an implementation-defect component). Its diagnosis and repair are Section~\ref{sec:arc}.

% ===========================================================================
